% SOURCE_AUTHORITY: sga5_fr_workpass.tex, lines 12650--12681 (LF-normalized unit).
% SOURCE_SHA256: 791F4EFFC5E02832D5D77ED03518C8156D6F07E4C8238B03545DB93D883FBB28
\subsubsection*{4.5}

Mais adiante necessitaremos de uma fórmula explícita para o caráter da representação de Swan. Com efeito, têm-se as fórmulas:
\[
\sigma^{*}_{y'}(s)=\sigma^{*}_{y''}(s)=\sigma_y(s)=
\begin{cases}
\displaystyle\sum_{\substack{p(z')=y\\s(z')=z'}}\bigl(1-v_{z'}(s)\bigr),& s\ne e,\\[0.8em]
\displaystyle\sum_{p(z')=y}\bigl(1+v_{z'}(\mathfrak D_{C'/C})\bigr)-n,& s=e,
\end{cases}
\tag{4.5.1}
\]
\[
\sigma^{\prime*}_{y'}(s)=\sigma^{\prime*}_{y''}(s)=\sigma'_y(s)=
\begin{cases}
\displaystyle\sum_{\substack{p(z')=y\\s(z')=z'}}\bigl(1-v_{z'}(s)\bigr),& s\ne e,\\[0.8em]
\displaystyle\sum_{p(z')=y}\bigl(1+v_{z'}(\mathfrak D_{C'/C})\bigr),& s=e,
\end{cases}
\tag{4.5.2}
\]
\[
a^{*}_{y'}(s)=a^{*}_{y''}(s)=a_y(s)=
\begin{cases}
\displaystyle-\sum_{\substack{p(z')=y\\s(z')=z'}}v_{z'}(s),& s\ne e,\\[0.8em]
\displaystyle\sum_{p(z')=y}v_{z'}(\mathfrak D_{C'/C}),& s=e.
\end{cases}
\tag{4.5.3}
\]
Para deduzir estas fórmulas, utiliza-se (3.3.1), tendo em conta que, fixado um ponto $y'_0\in C'$, para cada $y'\in C'$ tal que $p(y')=p(y'_0)$ pode-se encontrar um elemento $g'\in G$ tal que $g'(y')=y'_0$. Os elementos $g'$ constituem um sistema de representantes das classes à direita de $G_{y'_0}$ em $G$ e, além disso, tem-se
\[
G_{y'_0}=g'G_{y'}g'^{-1}.
\]
Este sistema de representantes pode ser utilizado no segundo membro da fórmula (3.3.1), e assim se obtêm as fórmulas escritas.
